\documentclass[Orbiter User Manual.tex]{subfiles}
\begin{document}

\section{Installation}
\label{sec:installation}
This section lists the computer requirements for running Orbiter on Linux, and contains download and installation instructions.

\subsection{Computer requirements}
\begin{sloppypar}
Orbiter runs on 64-bit (x86-64) Linux distributions with glibc 2.34 or newer and a Wayland or X11 desktop: Ubuntu 22.04 or newer and its relatives (Linux Mint 21, Pop!\_OS 22.04, Zorin OS 17 or newer), Debian 12 or newer, Fedora 36 or newer, RHEL, AlmaLinux and Rocky Linux 9 or newer, openSUSE Tumbleweed, Arch Linux, Manjaro, EndeavourOS and SteamOS 3. The release package includes Qt 6.8 and the other libraries that differ between distributions. The C library, the graphics drivers, X11 and Wayland, fonts, D-Bus, GLib, OpenSSL and PipeWire come from your distribution; the OpenOrbiter launcher checks them and offers to install any that are missing.\\
Orbiter needs a graphics driver with Vulkan 1.3 support: Mesa (AMD, Intel) or the NVIDIA driver. The VulkanClient uses the extensions VK\_\allowbreak KHR\_\allowbreak push\_\allowbreak descriptor, VK\_\allowbreak EXT\_\allowbreak shader\_\allowbreak object, VK\_\allowbreak EXT\_\allowbreak extended\_\allowbreak dynamic\_\allowbreak state3 and VK\_\allowbreak EXT\_\allowbreak vertex\_\allowbreak input\_\allowbreak dynamic\_\allowbreak state. Drivers without VK\_\allowbreak EXT\_\allowbreak shader\_\allowbreak object (Intel's ANV, older AMD Mesa) use the Khronos shader-object layer that comes with the package in lib/vulkan; the VulkanClient's log line \textit{Shader-object layer available} in Orbiter.log shows that it was found (see Doc/VulkanClient/VulkanClient.html). Sound is played through PipeWire; without it Orbiter runs silent.\\
\end{sloppypar}
\noindent
The following table lists the hardware requirements needed by Orbiter.

%\begin{table}[H]
	%\centering
	\begin{longtable}{ |p{0.065\textwidth}|p{0.3\textwidth}|p{0.55\textwidth}| }
	\hline\rule{0pt}{2ex}
	& \textbf{Minimum requirements} & \textbf{Recommended requirements}\\
	\hline\rule{0pt}{2ex}
	RAM & 500 MB & 2 GB\\
	\hline\rule{0pt}{2ex}
	CPU & Dual Core &\\
	\hline\rule{0pt}{2ex}
	GPU & 50 GFlops, Vulkan 1.3 driver & 100 GFlops\\
	\hline\rule{0pt}{2ex}
	Disk & 5 GB of free space & 10 GB of free space (80 GB if you want hi-res textures)\\
	\hline
	\end{longtable}
%\end{table}

\noindent
Orbiter supports an optional joystick.\\
The TrackIR head tracker is not available on Linux: its driver and the NPClient library Orbiter's TrackIR module uses exist only for Windows.

\subsection{Download}
The latest release of Orbiter 2027 for Linux can be downloaded from \url{https://github.com/racerx2/orbiter-2027-linux/releases}. The package is the file named Orbiter2027-Linux-x86\_64-<version>.tar.xz. Its SHA256 checksum is in the file of the same name with .sha256 added; \texttt{sha256sum -c Orbiter2027-Linux-x86\_64-<version>.tar.xz.sha256} in the download folder checks the download.\\
The surface textures of the planets and moons are a separate download, which the Launchpad offers when you launch a scenario (see below).\\
To build Orbiter yourself, follow COMPILE.md in the source repository. A package you build that way takes Qt (6.5 or newer) and the other libraries from your own distribution, and runs on distributions at least as new as the one it was built on.\\
If desired, high-resolution texture packs are available for Mercury, Earth, Moon, Mars, Titan and some minor-bodies, which can be found in the add-on repository of the Orbiter forum, \url{https://www.orbiter-forum.com/resources/}. The same packs are used on Windows and Linux. The archive files of a pack (Surf.tree, Elev.tree, Mask.tree and so on) belong in the Textures/<planet>/Archive folder of your Orbiter installation, for example Textures/Earth/Archive/Surf.tree.

\subsection{Installation}
\begin{itemize}
\item Unpack the package into a folder you own, for example your home folder, using either your file manager's archive tool or the command\\
\texttt{tar -xf Orbiter2027-Linux-x86\_64-<version>.tar.xz}\\
This creates the folder Orbiter2027-Linux-x86\_64-<version>, which is your Orbiter folder; you may rename or move it. Do not unpack it into /opt or another system folder: Orbiter writes its configuration, log and scenario files and the downloaded textures into its own folder.
\item If a previous version of Orbiter is already installed on your computer, you should not install the new version into the same folder, because this could lead to file conflicts. You may want to keep your old installation until you have made sure that the latest version works without problems. Multiple Orbiter installations can exist on the same computer.
\item After unpacking the package, make sure your Orbiter folder contains the executable (Orbiter), its launcher (OpenOrbiter), the lib folder with the included libraries and, among other files, the Config, Meshes, Scenarios and Textures subfolders.
\item Orbiter is launched with the OpenOrbiter launcher in the Orbiter folder, which starts the Orbiter executable from its own folder. The first time, the launcher checks the system libraries Orbiter uses, offers to install anything missing from your distribution's own package repositories (apt, dnf, zypper or pacman), and adds Orbiter to your application menu, so after that it starts from the menu like any other program. Orbiter's output goes to \textasciitilde/.cache/orbiter-linux/orbiter.log; \texttt{./OpenOrbiter -{}-verbose} keeps it in the terminal instead.
\item Graphics come from the included VulkanClient plugin, which renders with Vulkan: select it as graphics engine on the \textit{Video} tab of the Launchpad (see section \ref{ssec:launchpad_video}). Without a graphics engine, Orbiter runs in console mode, controlled from a text console (the terminal it was started from, or a console window of its own). Once running, Orbiter will show you the "Launchpad" dialog, where you can select video options and simulation parameters.
\item You are now ready to start, select a scenario from the Launchpad dialog and click the \textit{Launch Orbiter} button!
\item The surface textures of the planets and moons are a separate download. If some are missing when you launch a scenario, the Launchpad lists them with their sizes and offers to download and install them from \url{https://github.com/racerx2/Orbiter-linux-orbital-bodies-textures} (about 2.3 GB for all of them, 5 GB on disk once installed). Untick the bodies you don't want; \textit{Not now} launches without them, and \textit{Don't ask again} turns the offer off (PlanetTexCheck in Orbiter.cfg).
\end{itemize}

\noindent
\alertbox{If Orbiter does not show any scenarios in the Scenario tab, or if planets appear plain white without any textures when running the simulation, the most likely reason is that the installation packages were not properly unpacked. Make sure your Orbiter folder contains the subfolders as described above. If necessary, you may have to repeat the installation process.}

\subsection{Add-ons from the Windows version}
Orbiter on Linux looks up files the way Windows does: the case of file and folder names does not matter, and backslashes in paths inside configuration and scenario files are accepted. Add-ons made of meshes, textures, configuration files and scenarios therefore work without changes.\\
Add-ons that contain program modules (.dll files) are different: a Windows .dll cannot be loaded by the Linux version. Such an add-on needs a Linux build of its modules (.so files), made from its source code (see the Orbiter Developer Manual).\\
Help files (.chm) of Windows add-ons are Microsoft compiled HTML files, which Orbiter's help window on Linux cannot show (Orbiter's own .chm files on Linux are zip archives of HTML pages). Such help files can be read with a CHM viewer, for example xCHM or KchmViewer.

\subsection{Uninstall}
Run \texttt{./OpenOrbiter -{}-remove-menu} in the Orbiter folder to take Orbiter out of the application menu, then simply remove the Orbiter folders with all contents and subdirectories. This will completely remove Orbiter from your hard drive. The launcher's own files are in \textasciitilde/.cache/orbiter-linux (\textasciitilde/.cache/orbiter64-linux for the first Linux release) and can be removed as well. \texttt{-{}-remove-menu} keeps every copy of Orbiter out of the application menu; \texttt{./OpenOrbiter -{}-add-menu} puts it back.

\end{document}
